% Options for packages loaded elsewhere
\PassOptionsToPackage{unicode}{hyperref}
\PassOptionsToPackage{hyphens}{url}
\documentclass[
  ignorenonframetext,
]{beamer}
\newif\ifbibliography
\usepackage{pgfpages}
\setbeamertemplate{caption}[numbered]
\setbeamertemplate{caption label separator}{: }
\setbeamercolor{caption name}{fg=normal text.fg}
\beamertemplatenavigationsymbolsempty
% remove section numbering
\setbeamertemplate{part page}{
  \centering
  \begin{beamercolorbox}[sep=16pt,center]{part title}
    \usebeamerfont{part title}\insertpart\par
  \end{beamercolorbox}
}
\setbeamertemplate{section page}{
  \centering
  \begin{beamercolorbox}[sep=12pt,center]{section title}
    \usebeamerfont{section title}\insertsection\par
  \end{beamercolorbox}
}
\setbeamertemplate{subsection page}{
  \centering
  \begin{beamercolorbox}[sep=8pt,center]{subsection title}
    \usebeamerfont{subsection title}\insertsubsection\par
  \end{beamercolorbox}
}
% Prevent slide breaks in the middle of a paragraph
\widowpenalties 1 10000
\raggedbottom
\AtBeginPart{
  \frame{\partpage}
}
\AtBeginSection{
  \ifbibliography
  \else
    \frame{\sectionpage}
  \fi
}
\AtBeginSubsection{
  \frame{\subsectionpage}
}
\usepackage{iftex}
\ifPDFTeX
  \usepackage[T1]{fontenc}
  \usepackage[utf8]{inputenc}
  \usepackage{textcomp} % provide euro and other symbols
\else % if luatex or xetex
  \usepackage{unicode-math} % this also loads fontspec
  \defaultfontfeatures{Scale=MatchLowercase}
  \defaultfontfeatures[\rmfamily]{Ligatures=TeX,Scale=1}
\fi
\usepackage{lmodern}
\ifPDFTeX\else
  % xetex/luatex font selection
\fi
% Use upquote if available, for straight quotes in verbatim environments
\IfFileExists{upquote.sty}{\usepackage{upquote}}{}
\IfFileExists{microtype.sty}{% use microtype if available
  \usepackage[]{microtype}
  \UseMicrotypeSet[protrusion]{basicmath} % disable protrusion for tt fonts
}{}
\makeatletter
\@ifundefined{KOMAClassName}{% if non-KOMA class
  \IfFileExists{parskip.sty}{%
    \usepackage{parskip}
  }{% else
    \setlength{\parindent}{0pt}
    \setlength{\parskip}{6pt plus 2pt minus 1pt}}
}{% if KOMA class
  \KOMAoptions{parskip=half}}
\makeatother
\setlength{\emergencystretch}{3em} % prevent overfull lines
\providecommand{\tightlist}{%
  \setlength{\itemsep}{0pt}\setlength{\parskip}{0pt}}
% Auto-generated by infrastructure/rendering/_pdf_latex_helpers.py::extract_math_font_preamble.
% Minimal header passed to Pandoc via -H so Beamer slides pick up the same math font
% as the combined-PDF path. Adjust the manuscript's preamble.md to change the font globally.
\usepackage{unicode-math}
\setmathfont{latinmodern-math.otf}

% Slide layout defaults for warning-clean scientific decks.
\usepackage{etoolbox}
\IfFileExists{xurl.sty}{\usepackage{xurl}}{}
\IfFileExists{seqsplit.sty}{\usepackage{seqsplit}}{\newcommand{\seqsplit}[1]{#1}}
\protected\def\breakseq#1{\seqsplit{#1}}
\protected\def\breaktt#1{\begingroup\ttfamily\seqsplit{#1}\endgroup}
\setlength{\emergencystretch}{6em}
\tolerance=5000
\hbadness=10000
\hfuzz=1pt
\setlength{\tabcolsep}{2pt}
\AtBeginEnvironment{longtable}{\tiny\renewcommand{\arraystretch}{0.86}\setlength{\tabcolsep}{1pt}}
\AtBeginEnvironment{tabular}{\tiny\renewcommand{\arraystretch}{0.86}\setlength{\tabcolsep}{1pt}}

% Natbib and cross-reference fallbacks — slides don't load natbib
% or cleveref, but combined-PDF manuscript prose may emit these
% commands. The fallback renders citations as a bracketed key list
% and cross-references as detokenized labels so slides don't fail on
% undefined control sequences or raw underscores. \providecommand is
% a no-op if packages load later.
\providecommand{\citep}[1]{[#1]}
\providecommand{\citet}[1]{#1}
\providecommand{\citealp}[1]{#1}
\providecommand{\citeauthor}[1]{#1}
\providecommand{\citeyear}[1]{#1}
\providecommand{\cref}[1]{\texttt{\detokenize{#1}}}
\providecommand{\Cref}[1]{\texttt{\detokenize{#1}}}
\usepackage{bookmark}
\IfFileExists{xurl.sty}{\usepackage{xurl}}{} % add URL line breaks if available
\urlstyle{same}
% fallback for those not using the hyperref driver hyperxmp:
\makeatletter
\@ifundefined{xmpquote}{\newcommand{\xmpquote}[1]{#1}}{}
\makeatother
\hypersetup{
  hidelinks,
  pdfcreator={LaTeX via pandoc}}

\author{\texorpdfstring{}{}}
\date{}

\begin{document}

\section{Conclusion}\label{sec:conclusion}

\begin{frame}[allowframebreaks]{Conclusion}
This paper presented a small, tested domain language for specifying
controlled methods, informed by BPL's \breakseq{[@bpl2026]} domain-language
design for laboratory protocols and generalized to any controlled
procedure. It validates a simple proposition: a controlled vocabulary
expressed as typed, validated dataclasses --- not a parsed text grammar
--- is sufficient to reproduce BPL's core safety properties at a scope
appropriate for a template exemplar.
\end{frame}

\begin{frame}[fragile,allowframebreaks]{Exemplar achievements}
\protect\phantomsection\label{exemplar-achievements}
Operating as the methods-paper exemplar for the Research Project
Template methodology, the project deployed the three foundational
pillars:

\begin{enumerate}
\tightlist
\item
  \textbf{\breaktt{src/methods\_dsl/} library}: a controlled vocabulary,
  a dimensional unit system, four staged validation gates, a
  deterministic compiler, and four export formats --- with no plotting,
  no file I/O, and (with one declared logging exception) no
  \texttt{infrastructure} imports.
\item
  \textbf{\texttt{tests/} integrity}: a zero-mock suite over real
  \texttt{Method} fixtures covering the success path and every
  gate-failure mode, under a ≥90\% project coverage gate.
\item
  \textbf{\texttt{docs/} knowledge operations}: the correspondence with
  BPL's pipeline, testing philosophy, and operational rules that keep
  the library, scripts, and manuscript aligned.
\end{enumerate}
\end{frame}

\begin{frame}[fragile,allowframebreaks]{Technical contributions}
\protect\phantomsection\label{technical-contributions}
\begin{block}{Controlled vocabulary in code, not a parser}
\protect\phantomsection\label{controlled-vocabulary-in-code-not-a-parser}
The hallmark of this exemplar is the design choice it demonstrates: a
\texttt{.bpl} file's text grammar buys generality this template does not
need, while a frozen-dataclass model buys construction-time validation
(\texttt{\_\_post\_init\_\_}) that a parsed AST would have to re-derive.
The controlled vocabulary's discipline lives in the \emph{types}, not in
concrete syntax.
\end{block}

\begin{block}{Honest handling of trust boundaries}
\protect\phantomsection\label{honest-handling-of-trust-boundaries}
\texttt{trust.py}'s hash-chain documents its own limit explicitly: it
detects in-chain tampering but cannot detect a chain rewritten from
record zero. This makes the provenance guarantee a visible, testable
property rather than an implied cryptographic guarantee the
implementation does not actually provide.
\end{block}
\end{frame}

\begin{frame}[fragile,allowframebreaks]{Key insights}
\protect\phantomsection\label{key-insights}
\begin{enumerate}
\tightlist
\item
  \textbf{Determinism follows from explicit tie-breaking}: Kahn's
  algorithm \breakseq{[@kahn1962topological]} alone does not guarantee a stable
  schedule across runs --- the ascending-\texttt{step\_id} tie-break is
  what makes \texttt{plan\_hash} reproducible, and \breakseq{[@sec:results]}
  checks this live rather than asserting it.
\item
  \textbf{Staged short-circuit avoids misleading errors}: running
  \texttt{plan\_gate} and \texttt{target\_gate} against a structurally
  invalid method would report noise, not signal ---
  \texttt{run\_all\_gates} short-circuits after the first two gates for
  exactly this reason.
\item
  \textbf{A controlled vocabulary generalizes by restraint, not by
  expansion}: \breaktt{SensorCalibrationSweep} reuses every
  \texttt{StepKind} and \texttt{Target} the wet-lab-flavored
  \texttt{PBSPreparation} example uses; nothing was added to support a
  second domain.
\end{enumerate}
\end{frame}

\begin{frame}[fragile,allowframebreaks]{Future extensions}
\protect\phantomsection\label{future-extensions}
This foundation could be extended to:

\begin{itemize}
\tightlist
\item
  \textbf{A real \texttt{.bpl}-style parser}: add
  \texttt{grammar/}/\texttt{parser/}/\texttt{transformer/} stages ahead
  of the existing \texttt{model.py}, reusing every downstream gate and
  the compiler unchanged.
\item
  \textbf{More backends}: a \texttt{robot} execution target with
  capability-aware lowering, mirroring BPL's Biomek translation layer.
\item
  \textbf{A capability registry}: per-target primitive support
  declarations, generalizing BPL's \texttt{capabilities/} registry and
  \breaktt{bplc\ capabilities} report.
\end{itemize}
\end{frame}

\begin{frame}[fragile,allowframebreaks]{Final assessment}
\protect\phantomsection\label{final-assessment}
The \breaktt{template\_methods\_paper} tree is the canonical reference
for how a methods paper --- a manuscript whose subject is a methodology
--- stays synchronized with the code implementing that methodology
across rebuilds. The pipeline compiled both worked example methods,
wrote \breaktt{output/data/compiled\_plans.json},
\breaktt{output/reports/gate\_report.json}, and
\breaktt{output/reports/trust\_chain\_report.json}, and rendered this
markdown together with \texttt{config.yaml} into PDF.
\end{frame}

\end{document}
